\section*{Capítulo \ref{ch:g12:subduction-magmatism} --- O magmatismo das zonas de subducção}

\begin{solution}{exo:g12:subduction-magmatism:1}
(a) um \omterm{def:g12:subduction-magmatism:arc}{arco vulcânico}; (b) a \omterm{def:g12:subduction-magmatism:hydrous-melting}{fusão hidratada} (fusão por fluxo); (c) um
\omterm{def:g12:subduction-magmatism:granodiorite}{granodiorito}; (d) um \omterm{def:g12:subduction-magmatism:batholith}{batólito}; (e) o \omterm{def:g12:subduction-magmatism:continental-growth}{crescimento continental}.
\end{solution}

\begin{solution}{exo:g12:subduction-magmatism:2}
(1) Pelo menos $24/0.25 = \qty{96}{km/h}$, isto é, $\num{96000}/3600 \approx
\qty{27}{m/s}$. (2) A explosão foi pelo menos $480/96 = 5$ vezes mais rápida.
\end{solution}

\begin{solution}{exo:g12:subduction-magmatism:3}
(1) $838/1214 \approx 0.690$, cerca de \qty{69}{\percent}. (2) Do lado da
placa que se sobrepõe, a que não afunda. (3) Entre a fossa e a frente, a \omterm{def:g11:subduction-zones:slab}{laje}
ainda não está funda o bastante: ela não atingiu a profundidade, cerca de
cem quilômetros, em que a sua água é liberada num manto quente o bastante
para se fundir uma vez molhado.
\end{solution}

\begin{solution}{exo:g12:subduction-magmatism:4}
(1) $45/(45 + 15) = 0.75$: \qty{75}{\percent} do feldspato é plagioclásio.
(2) Sim: ela tem grãos grossos, contém mais de \qty{20}{\percent} de
quartzo, e a sua fração de plagioclásio fica entre 65 e \qty{90}{\percent}.
\end{solution}

\begin{solution}{exo:g12:subduction-magmatism:5}
Os magmas de arco são mais ricos em sílica, logo mais viscosos; e carregam a
água que os formou, logo são ricos em gás. Uma alta viscosidade e um alto
teor de gás favorecem ambos as erupções explosivas.
\end{solution}

\begin{solution}{exo:g12:subduction-magmatism:6}
(1) $H = 150 \times \tan\ang{35} \approx 150 \times 0.700 \approx
\qty{105}{km}$. (2) É a média mundial.
\end{solution}

\begin{solution}{exo:g12:subduction-magmatism:7}
$d = H/\tan\theta$: $105/\tan\ang{20} \approx 105/0.364 \approx
\qty{288}{km}$ a \ang{20}, e $105/\tan\ang{45} = \qty{105}{km}$ a
\ang{45}. Quanto mais suave a \omterm{def:g11:subduction-zones:slab}{laje}, mais longe da fossa fica o arco.
\end{solution}

\begin{solution}{exo:g12:subduction-magmatism:8}
(1) $T = -5.104 \times 3.16^2 + 132.9 \times 3.16 + 1120.7 \approx -51.0 +
420.0 + 1120.7 \approx \qty{1490}{\celsius}$. (2) Em cerca de $1490 - 1000 =
\qty{490}{\celsius}$. (3) Seca, não: \qty{1300}{\celsius} fica abaixo de
\qty{1490}{\celsius}. Molhada, sim: fica acima de \qty{1000}{\celsius}.
\end{solution}

\begin{solution}{exo:g12:subduction-magmatism:9}
(1) $2.5/2.7 \approx 0.93$ de um ano, cerca de 340 dias. (2) A avalanche
deslocou rocha antiga que já existia; os arcos produzem rocha nova a partir
do manto. Só os volumes são comparados.
\end{solution}

\begin{solution}{exo:g12:subduction-magmatism:10}
(1) $2.7 \times \num{1000000} = \qty{2.7e6}{km^3}$. (2) $\num{9e9}/2.7
\approx \num{3.3e9}$ anos, 3.3 bilhões de anos. (3) É menos que a idade da
Terra, então os arcos poderiam ter produzido toda a \omterm{def:g12:continental-crust:continental-crust}{crosta continental}; mas a
\omterm{def:g11:subduction-zones:subduction}{subducção} também leva crosta de volta, então, na verdade, grande parte da
crosta foi produzida e destruída mais de uma vez.
\end{solution}

\begin{solution}{exo:g12:subduction-magmatism:11}
(1) Os terremotos rasos também ocorrem na placa que se sobrepõe e na falha
entre as duas placas; só os profundos ficam apenas na \omterm{def:g11:subduction-zones:slab}{laje}.
(2) $116 - 80 = \qty{36}{km}$ em \qty{86.4}{km}: $\tan\theta = 36/86.4
\approx 0.417$, $\theta \approx \ang{22.6}$.
\end{solution}

\begin{solution}{exo:g12:subduction-magmatism:12}
No triângulo retângulo da proposição, a \omterm{def:g11:subduction-zones:slab}{laje} é a hipotenusa $L$ e $d$ o
cateto adjacente ao ângulo $\theta$, logo $\cos\theta = d/L$ e $L =
d/\cos\theta$. Para o Japão, $L = 270/\cos\ang{20} \approx 270/0.940 \approx
\qty{287}{km}$. A \qty{83}{mm} por ano: $\num{2.87e8}/83 \approx \num{3.5e6}$
anos, cerca de 3.5 milhões de anos.
\end{solution}

\begin{solution}{exo:g12:subduction-magmatism:13}
(1) $V = \pi r^2 h = \pi \times 0.5^2 \times 8 \approx \qty{6.3}{km^3}$.
(2) Cerca de $6.3/2.5 \approx 2.5$ vezes o volume da avalanche.
(3) $P = \rho g z$: a \qty{6}{km}, $2830 \times 9.81 \times 6000 \approx
\qty{1.67e8}{Pa} = \qty{0.17}{GPa}$; a \qty{14}{km}, cerca de
\qty{3.89e8}{Pa} = \qty{0.39}{GPa}.
\end{solution}

\begin{solution}{exo:g12:subduction-magmatism:14}
(1) $2.7 - 3.2 = \qty{-0.5}{km^3}$ por ano, uma perda de \qty{5e5}{km^3} por
milhão de anos. (2) $500 \times \num{5e5} = \qty{2.5e8}{km^3}$, isto é,
$\num{2.5e8}/\num{9e9} \approx \qty{2.8}{\percent}$ da crosta. (3) Cada taxa
é uma estimativa, com uma incerteza muito maior que a diferença de
\qty{0.5}{km^3} por ano entre elas: dentro dos erros, as duas se equilibram.
\end{solution}

\begin{solution}{exo:g12:subduction-magmatism:15}
(1) $\num{20000} \times 10 = \qty{2e5}{km^3}$. (2) $\num{2e5}/1 =
\num{200000}$ anos para todos os arcos continentais do mundo juntos. Um arco
recebe só uma pequena parte desse magma, então um único \omterm{def:g12:subduction-magmatism:batholith}{batólito} leva muito
mais tempo: ele é construído \omterm{def:g11:igneous-rocks:intrusion}{plúton} por \omterm{def:g11:igneous-rocks:intrusion}{plúton}, ao longo de milhões de anos.
(3) Pelo menos os \qty{5}{km} de rocha acima do seu topo.
\end{solution}

\begin{solution}{pb:g12:subduction-magmatism:1}
\textbf{1.} $111.2 \times \cos\ang{39} \approx 111.2 \times 0.777 \approx
\qty{86.4}{km}$.

\textbf{2.} $(144.1 - 141.0) \times 86.4 \approx 3.1 \times 86.4 \approx
\qty{268}{km}$, cerca de \qty{270}{km}.

\textbf{3.} $(141.0 - 139.7) \times 86.4 \approx \qty{112}{km}$.

\textbf{4.} Na placa que se sobrepõe, a que não afunda, a oeste da fossa: a
placa do Pacífico afunda sob ela.

\textbf{5.} Ali, a \omterm{def:g11:subduction-zones:slab}{laje} ainda não está funda o bastante para liberar a sua
água num manto quente o bastante para se fundir.

\textbf{6.} $116 - 80 = \qty{36}{km}$.

\textbf{7.} Um grau de longitude, \qty{86.4}{km}.

\textbf{8.} $\tan\theta = 36/86.4 \approx 0.417$, logo $\theta \approx
\ang{22.6}$.

\textbf{9.} A reta sobe \qty{36}{km} a cada \qty{86.4}{km} para leste, então
ela atinge a superfície $98/0.417 \approx \qty{235}{km}$ a leste de Iwate, a
$141.0 + 235/86.4 \approx \ang{143.7}$~Leste: cerca de $(144.1 - 143.7) \times 86.4
\approx \qty{35}{km}$ a oeste da fossa.

\textbf{10.} A \omterm{def:g11:subduction-zones:slab}{laje} não é um plano reto: ela entra na fossa quase
horizontal e se curva cada vez mais para baixo à medida que afunda.

\textbf{11.} $268 \times \tan\ang{20} \approx 268 \times 0.364 \approx
\qty{98}{km}$.

\textbf{12.} Eles concordam: \qty{98}{km} pelos dois caminhos.

\textbf{13.} Ambos ficam dentro da faixa mundial, um pouco mais rasos que a
média de \qty{105}{km}.

\textbf{14.} Megata fica \qty{112}{km} a oeste de Iwate, então a \omterm{def:g11:subduction-zones:slab}{laje} sob
ela está a cerca de $98 + 0.417 \times 112 \approx \qty{145}{km}$ de
profundidade.

\textbf{15.} No manto, cerca de $98 - 35 = \qty{63}{km}$ abaixo da base da
crosta.

\textbf{16.} Ela entrou perto da dorsal onde a crosta se formou, vinda da
água do mar; grande parte dela fica retida na serpentina.

\textbf{17.} $T = -5.104 \times 3.16^2 + 132.9 \times 3.16 + 1120.7 \approx
\qty{1490}{\celsius}$.

\textbf{18.} Em $1490 - 1000 = 490$ a $1490 - 950 = \qty{540}{\celsius}$:
cerca de \qty{500}{\celsius}.

\textbf{19.} A \omterm{def:g11:subduction-zones:slab}{laje} é várias centenas de graus mais fria que o manto em
volta dela, fria demais para se fundir mesmo com água: é a cunha do manto
quente acima dela, uma vez molhada, que se funde.

\textbf{20.} A \omterm{def:g11:subduction-zones:slab}{laje} fica a cerca de \textbf{\qty{100}{km} sob a frente
vulcânica} do nordeste do Japão (\qty{98}{km} pelos dois métodos).
\end{solution}
